%==============================================================================
% PARTIDA 01: CONEXION A LA RED EXTERNA DE MEDIDORES
% Codigo: OE.5.1.1.1
%==============================================================================
\subsection{OE.5.1.1.1 Acometida e Instalacion de Caja Portamedidor Monofasico}

\subsubsection{Especificaciones Tecnicas}

La acometida comprende el conjunto de elementos comprendidos entre la red publica de distribucion en baja tension y el medidor de energia electrica, incluyendo el cableado de alimentacion, la caja portamedidor monofasica y todos sus accesorios. La instalacion debe cumplir estrictamente con lo establecido en:

\begin{itemize}
  \item \textbf{CNE-U -- Seccion 2:} Reglas para instalaciones de acometida, Articulos 200.1 al 200.12 (conductores de acometida, soportes, altura libre, distancias de seguridad).
  \item \textbf{CNE-U -- Seccion 3:} Reglas para instalaciones de medidores, Articulos 300.1 al 300.8 (ubicacion, accesibilidad, proteccion contra intemperie).
  \item \textbf{RNE EM.010:} Articulo 8 (Acometidas) y Articulo 9 (Sistemas de medicion).
  \item \textbf{NTP 370.303:} Cajas portamedidores para energia electrica -- Requisitos de seguridad y desempeno.
  \item \textbf{NTP-IEC 60364-5-52:} Seleccion e instalacion de canalizaciones y conductores.
\end{itemize}

\paragraph{Caracteristicas de la caja portamedidor:}
\begin{itemize}[nosep]
  \item Tipo: Monofasico, para medidor de energia activa de 2 hilos (fase + neutro).
  \item Material: Termoplastico de alto impacto (policarbonato con estabilizador UV) o metalico con pintura electrostática anticorrosiva.
  \item Grado de proteccion: IP-54 minimo (proteccion contra polvo y salpicaduras de agua) segun NTP-IEC 60529.
  \item Tapa: Precintable, con visor transparente de policarbonato para lectura del medidor.
  \item Capacidad: Para alojar medidor monofasico de hasta 60 A y bornes de conexion de conductores hasta 16 mm2.
  \item Temperatura de operacion: -10 °C a +60 °C.
\end{itemize}

\paragraph{Conductor de acometida:}
\begin{itemize}[nosep]
  \item Tipo: Cable unipolar de cobre electrolitico recocido, clase 2 (cableado compacto).
  \item Aislamiento: Compuesto termoestable libre de halogenos (LSOH -- Low Smoke Zero Halogen), tipo NH-80, para tension nominal 0.6/1 kV.
  \item Seccion: 10 mm2 para conductor de fase y neutro (capacidad de corriente en ducto: 52 A a 30 °C segun CNE-U Tabla 5.2).
  \item Color: Negro para fase, blanco para neutro.
  \item Norma: NTP 370.252 (Cables aislados con compuesto termoplastico y termoestable).
\end{itemize}

\paragraph{Capacidad de corriente:}
La demanda maxima del proyecto es 21.73 A. El conductor de 10 mm2 con aislamiento LSOH soporta hasta 52 A en tuberia empotrada (CNE-U Tabla 5.2 para 3 conductores cargados), lo que representa un factor de seguridad del 139\% y permite ampliaciones futuras sin reemplazar el alimentador.

\subsubsection{Requisitos de Materiales}

\begin{longtable}{L{5.5cm} L{2.5cm} L{6.5cm}}
\toprule
\textbf{Material} & \textbf{Cantidad} & \textbf{Especificacion tecnica} \\
\midrule
Caja portamedidor monofasica & 1 UND & Termoplastica IP-54, con visor y tapa precintable, para medidor hasta 60 A \\
Conductor LSOH 10 mm2 negro & 15 m & Cobre electrolitico, 0.6/1 kV, fase \\
Conductor LSOH 10 mm2 blanco & 15 m & Cobre electrolitico, 0.6/1 kV, neutro \\
Conductor cobre desnudo 6 mm2 & 10 m & Cobre temple blando, clase 2, para conexion a tierra \\
Tuberia PVC pesada (PV-P) 1'' & 10 m & NTP 399.006, espesor nominal 2.5 mm, color naranja \\
Curva PVC pesada 1'' 90° & 2 pza & Radio medio segun NTP 399.006 \\
Union universal PVC pesada 1'' & 3 pza & Con pegamento integrado \\
Conector split-bolt bronce 6-10 mm2 & 2 pza & Para conexion de conductores en caja portamedidor \\
Terminal de compression tipo ojo 10 mm2 & 4 pza | Para conexion a bornes del medidor \\
Pegamento PVC (cemento solvente) & 1/8 gln | Para tuberia PVC, secado rapido \\
Cinta aislante autofundente 3M 23 o similar & 1 pza | Para sellado de conexiones \\
Cinta aislante vinilica 3M Tempflex o similar & 2 pza | Para identificacion y proteccion \\
Abrazadera tipo "U" galvanizada 1'' & 6 pza | Para fijacion de tuberia a muro \\
Perno de anclaje con tarugo 1/4'' x 2'' & 4 pza | Para fijacion de caja portamedidor \\
Gel dielectrico (compuesto siliconico) & 100 g & Para proteccion de conexiones contra corrosion y humedad \\
Precinto de seguridad plastico & 2 pza | Para sellado de tapa de caja portamedidor \\
\bottomrule
\end{longtable}

\subsubsection{Metodo de Ejecucion}

\begin{enumerate}
  \item \textbf{Replanteo:} Coordinar con la empresa distribuidora (Electro Puno S.A.A.) la ubicacion del punto de acometida y las condiciones de suministro. Marcar en fachada la ubicacion de la caja portamedidor segun plano IE-01, respetando:
  \begin{itemize}
    \item Altura: 1.50 m a 1.80 m sobre el nivel de la vereda.
    \item Distancia minima a esquinas: 0.30 m.
    \item Distancia minima a tuberias de gas: 1.00 m.
    \item Distancia minima a ventanas: 0.50 m horizontal, 0.30 m vertical.
  \end{itemize}
  \item \textbf{Fijacion de caja portamedidor:}
  \begin{enumerate}
    \item Perforar el muro con taladro percutor y broca de 1/4'' en las cuatro esquinas de la caja.
    \item Insertar tarugos de expansion de 1/4''.
    \item Colocar la caja portamedidor y fijar con pernos de anclaje, verificando nivelacion con nivel de burbuja (tolerancia: ±2 mm/m).
    \item Ajustar los pernos con llave de torque a 8 Nm.
  \end{enumerate}
  \item \textbf{Instalacion de tuberia de acometida:}
  \begin{enumerate}
    \item Trazar la ruta de la tuberia PVC pesada de 1'' desde el punto de conexion a red hasta la caja portamedidor.
    \item Corte de tuberia con segueta (corte perpendicular, sin rebabas).
    \item Aplicar pegamento PVC en el extremo del tubo y en el interior de curvas/uniones.
    \item Ensamblar asegurando penetracion de 2/3 de la longitud del accesorio.
    \item Fijar con abrazaderas tipo "U" cada 0.80 m como maximo.
    \item Pendiente minima del 0.5\% hacia la red para evitar acumulacion de humedad.
  \end{enumerate}
  \item \textbf{Tendido de conductores:}
  \begin{enumerate}
    \item Pasar guia de alambre galvanizado calibre 14 a traves de la tuberia.
    \item Fijar los conductores de fase (negro) y neutro (blanco) de 10 mm2 a la guia.
    \item Halar suavemente los conductores, evitando tirones bruscos que puedan danar el aislamiento.
    \item Dejar 0.50 m de conductor en cada extremo para conexiones.
  \end{enumerate}
  \item \textbf{Conexiones electricas:}
  \begin{enumerate}
    \item Pelar el aislamiento 2.0 cm con pelacables ajustable, sin danar el cobre.
    \item Instalar terminal de compression tipo ojo en los extremos de los conductores.
    \item Conectar fase al borne de entrada del medidor (borne L) y neutro al borne N.
    \item Conectar los conductores de salida del medidor a los bornes de carga.
    \item Aplicar gel dielectrico en todas las conexiones expuestas.
    \item Torque de apriete: 12 Nm para bornes de medidor, 8 Nm para conectores split-bolt.
  \end{enumerate}
  \item \textbf{Conexion de puesta a tierra:}
  \begin{enumerate}
    \item Conectar conductor de cobre desnudo 6 mm2 al borne de tierra de la caja portamedidor.
    \item Tender el conductor hasta la barra de tierra del TG-01 protegiendo con tuberia PVC en tramos expuestos.
  \end{enumerate}
  \item \textbf{Sellado y rotulado:}
  \begin{enumerate}
    \item Sellar las entradas de la caja portamedidor con masilla selladora (putty).
    \item Rotular la caja con adhesivo indeleble: numero de suministro, nombre del propietario y direccion.
    \item Instalar precintos de seguridad en la tapa.
  \end{enumerate}
\end{enumerate}

\subsubsection{Control de Calidad}

\begin{longtable}{L{4.5cm} L{4.5cm} L{4.5cm}}
\toprule
\textbf{Parametro} & \textbf{Metodo de verificacion} & \textbf{Criterio de aceptacion} \\
\midrule
Continuidad electrica & Multimetro digital en modo continuidad (escala < 10 $\Omega$) & Resistencia < 1 $\Omega$ en conductores fase, neutro y tierra \\
Resistencia de aislamiento & Megohmetro 1000 V entre fase-neutro, fase-tierra, neutro-tierra & > 1 M$\Omega$ (recomendado > 10 M$\Omega$ para instalacion nueva) \\
Polaridad & Verificador de polaridad en bornes del medidor & Fase en borne L, neutro en borne N \\
Nivelacion de caja & Nivel de burbuja en dos ejes ortogonales & Tolerancia: ±2 mm/m \\
Fijacion de caja & Prueba de traccion manual (10 kg) & Sin desplazamiento ni vibracion \\
Estanqueidad & Inspeccion visual de sellos y empaquetaduras & Sin aberturas ni espacios visibles \\
Precintos & Inspeccion visual & Instalados y numerados \\
Torque de conexiones & Llave de torque dinamometrica & 12 Nm en bornes de medidor, 8 Nm en split-bolt \\
\bottomrule
\end{longtable}

\paragraph{Pruebas electricas obligatorias:}
\begin{enumerate}
  \item Prueba de continuidad de todos los conductores (fase, neutro, tierra).
  \item Prueba de resistencia de aislamiento con megohmetro de 1000 V (minimo 1 M$\Omega$).
  \item Prueba de polaridad en bornes de entrada y salida del medidor.
  \item Prueba de tension con multimetro (debe medir 220 V ± 10\% entre fase y neutro).
  \item Prueba de corriente de fuga con pinza amperimetrica de fuga (< 5 mA).
  \item Verificacion de funcionamiento del medidor (lectura inicial registrada).
\end{enumerate}

\subsubsection{Metodo de Medicion}

La medicion se efectuara por unidad (UND) de caja portamedidor monofasico completamente instalada, incluyendo la caja, conductores desde el punto de acometida hasta los bornes de carga del medidor, tuberia de proteccion, accesorios, conexiones, sellado, rotulado y pruebas electricas. Se considera una unidad por cada suministro electrico independiente.

\subsubsection{Unidad de Medida}

UND (Unidad) -- Una caja portamedidor monofasico instalada, conectada y operativa, con todos sus accesorios.

\subsubsection{Forma de Pago}

El pago se efectuara por unidad completamente instalada, previa verificacion de los siguientes hitos:
\begin{enumerate}
  \item Aprobacion del replanteo y ubicacion por parte de supervision.
  \item Instalacion completa y conexion de conductores.
  \item Pruebas de calidad aprobadas (continuidad, aislamiento, polaridad).
  \item Verificacion de funcionamiento del medidor por la empresa distribuidora.
\end{enumerate}

El precio unitario incluye: suministro de caja portamedidor, conductores, tuberia, accesorios, conectores, selladores, mano de obra calificada, herramientas, equipos de prueba y todos los recursos necesarios para dejar la partida en condiciones de operacion. No incluye el medidor de energia (suministrado por la empresa distribuidora).
